Télécom Paris, école de l’IMT (Institut Mines-Télécom) et membre fondateur de l’Institut Polytechnique de Paris, est une grande école du top 5 des écoles d’ingénieurs généralistes françaises. Institution à « taille humaine » mais à forte composante internationale, Télécom Paris est reconnue pour sa proximité avec les entreprises. Cette école publique garantit une excellente employabilité dans tous les secteurs et apparait comme la 1ère grande école d’ingénieurs sur toute la verticale du numérique (des couches matérielles jusqu’aux usages).

Avec des enseignements d’excellence et une pédagogie innovante, Télécom Paris est au cœur d'un écosystème d'innovation unique, fondé sur l’interaction et l’importance du mode projet dans sa formation d’une part, et de l’autre part sur sa recherche interdisciplinaire. Ses enseignants-chercheurs sont affiliés à deux laboratoires de recherche: d'une part, le laboratoire LTCI qui est présenté par le HCERES comme une unité phare dans le domaine des sciences du numérique avec un rayonnement remarquable à l’international; et d'autre part, le laboratoire i3, Institut interdisciplinaire de l’Innovation (I3 - UMR 9217 du CNRS), qui poursuit un programme de recherche multidisciplinaire centré sur l’innovation dans le cadre d'une collaboration avec l’École Polytechnique et Mines ParisTech.

Basée à Palaiseau, au cœur du campus de l’Institut Polytechnique aux côtés de l’École polytechnique, de l’ENSTA, de Télécom Sud Paris et de l’ENSAE, Télécom Paris est également dotée d’un incubateur basé à Paris au cœur de l’écosystème français des start-ups.

Télécom Paris se positionne comme le laboratoire à ciel ouvert de tous les grands défis technologiques et sociétaux : intelligence artificielle, informatique quantique, IOT, cybersécurité, grands équipements numériques (Cloud), 5G/6G, Green IT.